\documentclass[12pt,a4paper]{article}

% ─── Packages ────────────────────────────────────────────────────────────────
\usepackage[margin=2.5cm]{geometry}
\usepackage{graphicx}
\usepackage{booktabs}
\usepackage{amsmath}
\usepackage{hyperref}
\usepackage{listings}
\usepackage{xcolor}
\usepackage{float}
\usepackage{caption}
\usepackage{subcaption}
\usepackage{parskip}      % space between paragraphs, no indent
\usepackage{microtype}
\usepackage{lmodern}

% ─── Code listing style ───────────────────────────────────────────────────────
\lstset{
  basicstyle=\ttfamily\small,
  backgroundcolor=\color{gray!10},
  frame=single,
  breaklines=true,
  columns=fullflexible,
  keepspaces=true,
}

% ─── Hyperlink colours ────────────────────────────────────────────────────────
\hypersetup{
  colorlinks=true,
  linkcolor=blue!60!black,
  urlcolor=blue!60!black,
  citecolor=blue!60!black,
}

% ─── Title ───────────────────────────────────────────────────────────────────
\title{%
  \textbf{Assignment 2 --- IMDB Sentiment Classification with CNNs}\\[6pt]
  \large DATA-MSML 612 --- Deep Learning\\
  University of Maryland
}
\author{Shrikanth Vilvadrinath}
\date{March 6, 2026}

% ─────────────────────────────────────────────────────────────────────────────
\begin{document}
\maketitle
\tableofcontents
\newpage

% =============================================================================
\section{Introduction}
% =============================================================================

This report presents a comparative study of three model families for binary sentiment
classification on the IMDB movie-review dataset: Logistic Regression (LR), a Feed-Forward
Neural Network (FFNN), and a Convolutional Neural Network (CNN). We additionally
investigate the effect of stopword removal on CNN performance. The target accuracy
threshold throughout is \textbf{>85\%}.

The results reveal a nuanced performance landscape: the FFNN slightly outperforms the CNN
in test accuracy (86.50\% vs.\ 85.57\%), which is counter-intuitive given the CNN's richer
inductive bias. Both neural models exhibit significant overfitting, which the analysis
examines in depth. All three primary models clear the 85\% threshold.

\textbf{Framework:} TensorFlow 2.10 / Keras. \textbf{Seed:} 42 (fully reproducible).

% =============================================================================
\section{Dataset}
% =============================================================================

The IMDB dataset \cite{maas2011} contains \textbf{50,000 movie reviews} split evenly as
25,000 for training and 25,000 for testing. Labels are binary (1 = positive, 0 = negative).
The dataset is perfectly balanced: 12,500 positive and 12,500 negative examples in each
split.

\begin{quote}
\textit{Note on the assignment description:} The assignment text states ``50,000 reviews
for training and the same amount again for testing.'' The actual Keras IMDB dataset
provides \textbf{25,000 training and 25,000 test reviews (50,000 total)}. The numbers
used here reflect the actual dataset.
\end{quote}

\subsection*{Preprocessing}
\begin{itemize}
  \item Vocabulary limited to the top \textbf{5,000} most frequent words
        (\texttt{NUM\_WORDS=5000}; rarer words replaced with the \texttt{$<$UNK$>$} token).
  \item Reviews padded/truncated to \textbf{500 tokens} (\texttt{MAX\_LEN=500}) using
        pre-padding (zeros prepended to short reviews; long reviews truncated from the
        left).
  \item Raw review length statistics: min = 7, max = 2,494, mean = 234.8, median = 176.0.
\end{itemize}

\subsection*{Sample Decoded Review}
To illustrate the preprocessing pipeline, training review 0 (Positive label) decodes as:

\begin{quote}
\textit{``this film was just brilliant casting location scenery story direction
everyone's really suited the part they played and you could just imagine being there
robert redford's is an amazing actor\ldots''}
\end{quote}

Reviews are stored as integer word-rank sequences and decoded back via the IMDB word
index. The model sees integer tokens, not raw text.

% =============================================================================
\section{Models and Architecture}
% =============================================================================

\subsection{Logistic Regression (Baseline)}

\textbf{Input:} Multi-hot binary bag-of-words (BoW) vectors of shape
$(25{,}000,\ 5{,}000)$. Position $i$ is 1 if word-ID $i$ appears in the review, 0
otherwise.

Using raw padded integer sequences as LR input is semantically incorrect: word-rank 42
would be treated as twice as important as word-rank 21, which is meaningless. Multi-hot
BoW is the standard representation for linear text classification.

\begin{lstlisting}[caption={Bag-of-Words encoding}]
def to_multi_hot(sequences, num_words):
    result = np.zeros((len(sequences), num_words), dtype=np.float32)
    for i, seq in enumerate(sequences):
        indices = seq[seq > 0]              # skip padding
        indices = indices[indices < num_words]
        result[i, indices] = 1.0
    return result
\end{lstlisting}

\subsection{Feed-Forward Neural Network (Baseline)}

\begin{lstlisting}[caption={FFNN architecture}]
Embedding(5000, 32, input_length=500)  ->  (batch, 500, 32)
Flatten()                              ->  (batch, 16000)
Dense(250, activation='relu')          ->  (batch, 250)
Dense(1,   activation='sigmoid')       ->  (batch, 1)
# Total parameters: 4,160,501
\end{lstlisting}

The FFNN uses a \emph{learned} Embedding layer, giving it representation-learning
capability. This is standard practice for neural text classification and provides a
meaningful comparison point before adding convolutional layers. The model learns
distributed word representations jointly with the classification head.

\subsection{Convolutional Neural Network (CNN)}

\begin{lstlisting}[caption={CNN architecture with dimension walkthrough}]
Embedding(5000, 32, input_length=500)          -> (batch, 500, 32)
Conv1D(32, kernel_size=7, activation='relu')   -> (batch, 494, 32) [500-7+1=494]
MaxPooling1D(pool_size=5)                      -> (batch,  98, 32) [floor(494/5)=98]
Flatten()                                      -> (batch, 3136)    [98*32=3136]
Dense(250, activation='relu')                  -> (batch, 250)
Dense(1,   activation='sigmoid')               -> (batch, 1)
# Total trainable parameters: 951,701
\end{lstlisting}

\textbf{Architecture rationale:}
\begin{itemize}
  \item \texttt{kernel\_size=7}: Captures 7-word n-gram patterns. A window of 7 is large
        enough to cover sentiment negation phrases such as \textit{``not at all good''} or
        \textit{``could not have been better''}.
  \item \texttt{pool\_size=5}: Aggressive subsampling extracts the most salient feature
        activations, reducing the temporal dimension $494 \to 98$.
  \item \texttt{Dense(250)}: Standard-capacity classification head; 98\,$\times$\,32 = 3,136
        flattened inputs provide sufficient signal.
  \item \texttt{sigmoid} output with \texttt{binary\_crossentropy} loss for binary
        classification.
\end{itemize}

\textbf{Compile:} \texttt{optimizer='adam'}, \texttt{loss='binary\_crossentropy'},
\texttt{metrics=['accuracy']}.\\
\textbf{Training:} 10 epochs, \texttt{batch\_size=128}, \texttt{validation\_split=0.2}
(test set evaluated only once after training to avoid data leakage).

% =============================================================================
\section{Results}
% =============================================================================

All models trained with \texttt{SEED=42} (\texttt{random}, \texttt{numpy}, and
\texttt{tensorflow} seeds all set). Table~\ref{tab:results} summarises the outcomes.

\begin{table}[H]
  \centering
  \caption{Model accuracy results on IMDB dataset (SEED=42)}
  \label{tab:results}
  \begin{tabular}{lrrr}
    \toprule
    \textbf{Model} & \textbf{Train Acc} & \textbf{Test Acc} & \textbf{vs.\ 85\%} \\
    \midrule
    Logistic Regression (BoW)   & 95.82\% & 85.71\% & PASS +0.71\% \\
    Feed-Forward NN (Embedding) & 97.40\% & 86.50\% & PASS +1.50\% \\
    CNN                         & 97.34\% & 85.57\% & PASS +0.57\% \\
    CNN (No Stopwords)          & 97.09\% & 84.80\% & FAIL $-$0.20\% \\
    \bottomrule
  \end{tabular}
\end{table}

\textbf{All three primary models exceed the 85\% accuracy target.}

% =============================================================================
\section{Analysis}
% =============================================================================

\subsection{Model Performance Hierarchy}

Test accuracy ranking: \textbf{FFNN (86.50\%) $>$ LR (85.71\%) $>$ CNN (85.57\%)
$>$ CNN-NS (84.80\%)}.

The headline finding is that the FFNN outperforms the CNN. This is counter-intuitive given
the CNN's richer inductive bias, and is examined in Section~\ref{sec:ffnn-vs-cnn}.

\subsection{FFNN vs.\ CNN --- Why the Simpler Model Wins Here}
\label{sec:ffnn-vs-cnn}

The FFNN has \textbf{4.16 million parameters} versus the CNN's \textbf{951,701} ---
roughly 4$\times$ more capacity. With 25,000 training examples, the FFNN's
Embedding + Flatten path learns global word co-occurrence patterns across the entire
500-token sequence simultaneously.

\textbf{Training curve analysis} (Figures~\ref{fig:cnn-acc} and \ref{fig:cnn-loss}):

Both models overfit significantly. The FFNN exhibits more extreme overfitting:
\begin{itemize}
  \item Training accuracy climbs to $\approx$99.99\% by epoch 10.
  \item Validation accuracy plateaus at $\approx$85--86\%.
  \item Training loss collapses to 0.0004 while \texttt{val\_loss} rises to 0.9988
        (a 250$\times$ divergence).
\end{itemize}
The CNN overfits less severely (train 97.34\% vs.\ test 85.57\%, an 11.77\% gap), and
its validation loss curve shows divergence from training loss starting around epoch 2--3.

\textbf{Why CNNs typically dominate NLP:} Conv1D's advantage becomes pronounced with
(1) longer sequences, (2) pre-trained embeddings (GloVe, BERT), or (3) scarce training
data where parameter efficiency matters. For this specific dataset (25k samples,
vocabulary 5k, sequences padded to 500), the FFNN's global token-position attention
provides sufficient discriminative signal.

\textbf{Recommended fix:} Adding \texttt{Dropout(0.3--0.5)} after \texttt{Dense(250)},
or applying early stopping at epoch 5--7 based on validation loss, would reduce
overfitting and likely close the performance gap between models.

\subsection{Logistic Regression as a Strong Baseline}

LR achieves \textbf{85.71\% test accuracy} --- only 0.79\% behind the FFNN and 0.14\%
ahead of the CNN. This demonstrates that the IMDB task has strong lexical signal: explicit
sentiment vocabulary (``excellent'', ``terrible'', ``masterpiece'', ``boring'') that a
simple BoW representation can exploit without any learned representations.

The 10.11\% train--test gap (95.82\% vs.\ 85.71\%) reflects overfitting despite sklearn's
default L2 regularisation (\texttt{C=1.0}). Reducing \texttt{C} (e.g., \texttt{C=0.1})
would likely close this gap.

\subsection{Effect of Stopword Removal (Task 6)}

Removing NLTK English stopwords reduces CNN test accuracy from \textbf{85.57\% to
84.80\%}, a drop of \textbf{$-$0.77 percentage points}.

\subsubsection*{The Negation Story}

NLTK's English stopword list includes sentiment-critical negation words:

\begin{quote}
\textit{no, not, nor, couldn't, didn't, doesn't, hadn't,
hasn't, haven't, isn't, mustn't, needn't, shouldn't, wasn't, weren't, won't,
wouldn't} (17 of 20 common negation words; \textit{never}, \textit{neither},
and \textit{cannot} are not in NLTK's stopword list)
\end{quote}

These are among the most powerful sentiment modifiers in natural language.

\subsubsection*{Concrete Before/After Decode Example}

Table~\ref{tab:decode} shows a training review (review \#5, Negative label) before
and after stopword removal. Note how function words (``than'', ``it'', ``that'',
``the'', ``all'', ``other'') are stripped alongside the structurally important
``it's'' and ``those''.

\begin{table}[H]
  \centering
  \caption{Stopword removal: before and after decode of a sample review (\#5)}
  \label{tab:decode}
  \begin{tabular}{p{1.5cm} p{11cm}}
    \toprule
    \textbf{Stage} & \textbf{First 20 decoded words} \\
    \midrule
    Before &
      \texttt{begins better than it ends funny that the russian $<$UNK$>$ crew
      $<$UNK$>$ all other actors it's like those scenes} \\
    \addlinespace
    After &
      \texttt{begins better ends funny russian $<$UNK$>$ crew $<$UNK$>$ actors
      like scenes documentary shots br br spoiler part message $<$UNK$>$} \\
    \bottomrule
  \end{tabular}
\end{table}

When ``not'' precedes a positive phrase (\textit{``not a great film''} $\to$
\textit{``great film''}) or ``never'' precedes a negative one (\textit{``never
disappoints''} $\to$ \textit{``disappoints''}), the sentiment inverts. The CNN's 7-gram
Conv1D filters learn these negation n-gram patterns as fixed-window features. Removing
the negation token disrupts the n-gram and causes the filter to produce the wrong
activation.

The relatively small drop ($-$0.77\%) rather than a catastrophic collapse suggests:
\begin{enumerate}
  \item Many reviews use explicit sentiment vocabulary that does not rely on negation.
  \item The CNN partially recovers from surrounding context.
  \item Neutral stopwords (``the'', ``a'', ``is'') dominate the $\approx$111 removed
        words per review, diluting the impact of the 17 negation words.
\end{enumerate}

% =============================================================================
\section{Training Curves}
% =============================================================================

\begin{figure}[H]
  \centering
  \includegraphics[width=0.90\textwidth]{output/task5_cnn_accuracy.png}
  \caption{CNN training vs.\ validation accuracy per epoch. Both metrics are shown in
           percentage (\%). The gap between training ($\approx$97\%) and validation
           ($\approx$85\%) accuracy at epoch 10 indicates significant overfitting.}
  \label{fig:cnn-acc}
\end{figure}

\begin{figure}[H]
  \centering
  \includegraphics[width=0.90\textwidth]{output/task5_cnn_loss.png}
  \caption{CNN training vs.\ validation loss (binary cross-entropy) per epoch. Validation
           loss begins rising after epoch 2--3 while training loss continues to drop ---
           a classic overfitting signature.}
  \label{fig:cnn-loss}
\end{figure}

\begin{figure}[H]
  \centering
  \includegraphics[width=0.90\textwidth]{output/task7_model_comparison.png}
  \caption{Model comparison: grouped bar chart of training and test accuracy for all four
           models. The horizontal dashed line marks the 85\% target. Value labels show
           exact percentages. Note that all three primary models (LR, FFNN, CNN) clear the
           85\% threshold; the CNN without stopwords falls marginally short.}
  \label{fig:comparison}
\end{figure}

% =============================================================================
\section{Conclusion}
% =============================================================================

All three primary models exceed the 85\% accuracy target on the IMDB test set:
LR 85.71\%, FFNN 86.50\%, CNN 85.57\%. Key conclusions:

\begin{enumerate}
  \item \textbf{Bag-of-words LR is a competitive baseline} for sentiment
        classification --- it actually outperforms the CNN by 0.14 percentage points
        (85.71\% vs.\ 85.57\%), because the IMDB task has strong lexical signal that
        a linear model can exploit directly.

  \item \textbf{The FFNN outperforms the CNN} due to its 4$\times$ larger parameter
        count and sufficient training data for global pattern learning. Both models
        overfit heavily; Dropout or early stopping are the recommended remedies.

  \item \textbf{Stopword removal hurts CNN accuracy ($-$0.77\%)} because NLTK
        stopwords include 17 negation words that are critical for sentiment n-gram
        patterns (e.g., \textit{``not good''}, \textit{``never disappoints''}).
\end{enumerate}

% =============================================================================
\begin{thebibliography}{9}

\bibitem{maas2011}
  Maas, A.\,L., Daly, R.\,E., Pham, P.\,T., Huang, D., Ng, A.\,Y., and
  Potts, C. (2011). \textit{Learning word vectors for sentiment analysis.}
  Proceedings of ACL 2011, pp.\ 142--150.

\bibitem{chollet2018}
  Chollet, F. (2018). \textit{Deep Learning with Python.} Manning Publications.

\bibitem{keras}
  Keras documentation. \url{https://keras.io/api/layers/}

\bibitem{nltk}
  NLTK stopwords corpus. \url{https://www.nltk.org/}

\end{thebibliography}

\vspace{1em}
\noindent\small\textit{Code: \texttt{Assignment 2/code/assignment2.py} \quad
Output plots: \texttt{Assignment 2/output/}}

\end{document}
